\documentclass{controle}
\usepackage{main}

\title{Évaluation n°14 : Arithmétique}
\author{Seconde 3}
\date{23 Mars 2026}

% \printanswers

\begin{document}
\begin{pycode}
from random import *
from functools import reduce
from operator import mul
import math

def gen_decompose(s):
  seed(s)
  p2 = randrange(0,6)
  p3 = randrange(0,3)
  p5 = randrange(0,3)
  p7 = randrange(0,2)
  p11 = randrange(0,2)
  p13 = 1 if p7 == 0 and p11 == 0 else 0

  n = 2**p2 * 3**p3 * 5**p5 * 7**p7 * 11**p11 * 13**p13
  sol = [2]*p2 + [3]*p3 + [5]*p5 + [7]*p7 + [11]*p11 + [13]*p13
  return (n,sol)

def simplify_fracts(sol1,sol2):
  l1, l2 = len(sol1), len(sol2)
  c1, c2 = 0, 0
  cancel1, cancel2 = [], []
  without1, without2 = [], []
  while c1 < l1 or c2 < l2:
    elt1 = sol1[c1] if c1 < l1 else math.inf
    elt2 = sol2[c2] if c2 < l2 else math.inf
    if elt1 == elt2:
      cancel1.append(f"\\cancel{{{elt1}}}")
      cancel2.append(f"\\cancel{{{elt2}}}")
      c1 += 1
      c2 += 1
    elif c2 == l2 or elt1 < elt2:
      cancel1.append(f"{elt1}")
      without1.append(elt1)
      c1 += 1
    else:
      cancel2.append(f"{elt2}")
      without2.append(elt2)
      c2 += 1
  without1 = without1 if len(without1) > 0 else [1]
  without2 = without2 if len(without2) > 0 else [1]

  str_cancel1 = " \\times ".join(cancel1) 
  str_cancel2 = " \\times ".join(cancel2)
  str_without1 = " \\times ".join(map(str,without1)) 
  str_without2 = " \\times ".join(map(str,without2))
  res1 = reduce(mul, without1, 1)
  res2 = reduce(mul, without2, 1)


  return str_cancel1, str_cancel2, str_without1, str_without2, res1, res2, all(map(lambda x: x==2 or x==5,without2))
\end{pycode}

\titleandname{}
\version{}
\begin{questions}
\question
\begin{pycode}
n1,sol1 = gen_decompose(1)
n2,sol2 = gen_decompose(20)
n3,sol3 = gen_decompose(3)
n4,sol4 = gen_decompose(4)
n5,sol5 = gen_decompose(5)

#Fractions
numc12, denc12, num12, den12, numr12, denr12, deci12 = simplify_fracts(sol1,sol2)
numc34, denc34, num34, den34, numr34, denr34, deci34 = simplify_fracts(sol3,sol4)
numc52, denc52, num52, den52, numr52, denr52, deci52 = simplify_fracts(sol5,sol2)

\end{pycode}
\begin{parts}
\part Donner la décomposition en facteurs premiers des nombres suivants.
\begin{subparts}
\subpart $\py{n1}$
\subpart $\py{n2}$
\subpart $\py{n3}$
\subpart $\py{n4}$
\subpart $\py{n5}$
\end{subparts}

\vspace*{0.5cm}

\part En déduire l'écriture sous forme irréductible des fractions suivantes.
\begin{subparts}
\subpart $\dfrac{\py{n1}}{\py{n2}} = $
\subpart $\dfrac{\py{n3}}{\py{n4}} = $
\subpart $\dfrac{\py{n5}}{\py{n2}} = $
\end{subparts}

\vspace*{0.5cm}

\part En déduire pour chacune des fractions ci-dessus si elle correspond à un nombre décimal ou non.

\end{parts}

\begin{solution}
\begin{parts}
\part
\begin{subparts}
\subpart $\py{n1} = \py{" \\times ".join([str(elt) for elt in sol1])}$
\subpart $\py{n2} = \py{" \\times ".join([str(elt) for elt in sol2])}$
\subpart $\py{n3} = \py{" \\times ".join([str(elt) for elt in sol3])}$
\subpart $\py{n4} = \py{" \\times ".join([str(elt) for elt in sol4])}$
\subpart $\py{n5} = \py{" \\times ".join([str(elt) for elt in sol5])}$
\end{subparts}
\part
\begin{subparts}
\subpart $\dfrac{\py{n1}}{\py{n2}} = \dfrac{\py{numc12}}{\py{denc12}} = \dfrac{\py{num12}}{\py{den12}} = \dfrac{\py{numr12}}{\py{denr12}}$
\subpart $\dfrac{\py{n3}}{\py{n4}} = \dfrac{\py{numc34}}{\py{denc34}} = \dfrac{\py{num34}}{\py{den34}} = \dfrac{\py{numr34}}{\py{denr34}}$
\subpart $\dfrac{\py{n5}}{\py{n2}} = \dfrac{\py{numc52}}{\py{denc52}} = \dfrac{\py{num52}}{\py{den52}} = \dfrac{\py{numr52}}{\py{denr52}}$
\end{subparts}
\part
\begin{subparts}
\subpart \py{"Décimal" if deci12 else "Non Décimal"}
\subpart \py{"Décimal" if deci34 else "Non Décimal"}
\subpart \py{"Décimal" if deci52 else "Non Décimal"}
\end{subparts}
\end{parts}
\end{solution}
\end{questions}

\newpage
\titleandname{}
\version{}
\begin{questions}
\question
\begin{pycode}
n1,sol1 = gen_decompose(10)
n2,sol2 = gen_decompose(200)
n3,sol3 = gen_decompose(30)
n4,sol4 = gen_decompose(40)
n5,sol5 = gen_decompose(50)

#Fractions
numc12, denc12, num12, den12, numr12, denr12, deci12 = simplify_fracts(sol1,sol2)
numc34, denc34, num34, den34, numr34, denr34, deci34 = simplify_fracts(sol3,sol4)
numc52, denc52, num52, den52, numr52, denr52, deci52 = simplify_fracts(sol5,sol2)

\end{pycode}
\begin{parts}
\part Donner la décomposition en facteurs premiers des nombres suivants.
\begin{subparts}
\subpart $\py{n1}$
\subpart $\py{n2}$
\subpart $\py{n3}$
\subpart $\py{n4}$
\subpart $\py{n5}$
\end{subparts}

\vspace*{0.5cm}

\part En déduire l'écriture sous forme irréductible des fractions suivantes.
\begin{subparts}
\subpart $\dfrac{\py{n1}}{\py{n2}} = $
\subpart $\dfrac{\py{n3}}{\py{n4}} = $
\subpart $\dfrac{\py{n5}}{\py{n2}} = $
\end{subparts}

\vspace*{0.5cm}

\part En déduire pour chacune des fractions ci-dessus si elle correspond à un nombre décimal ou non.

\end{parts}

\begin{solution}
\begin{parts}
\part
\begin{subparts}
\subpart $\py{n1} = \py{" \\times ".join([str(elt) for elt in sol1])}$
\subpart $\py{n2} = \py{" \\times ".join([str(elt) for elt in sol2])}$
\subpart $\py{n3} = \py{" \\times ".join([str(elt) for elt in sol3])}$
\subpart $\py{n4} = \py{" \\times ".join([str(elt) for elt in sol4])}$
\subpart $\py{n5} = \py{" \\times ".join([str(elt) for elt in sol5])}$
\end{subparts}
\part
\begin{subparts}
\subpart $\dfrac{\py{n1}}{\py{n2}} = \dfrac{\py{numc12}}{\py{denc12}} = \dfrac{\py{num12}}{\py{den12}} = \dfrac{\py{numr12}}{\py{denr12}}$
\subpart $\dfrac{\py{n3}}{\py{n4}} = \dfrac{\py{numc34}}{\py{denc34}} = \dfrac{\py{num34}}{\py{den34}} = \dfrac{\py{numr34}}{\py{denr34}}$
\subpart $\dfrac{\py{n5}}{\py{n2}} = \dfrac{\py{numc52}}{\py{denc52}} = \dfrac{\py{num52}}{\py{den52}} = \dfrac{\py{numr52}}{\py{denr52}}$
\end{subparts}
\part
\begin{subparts}
\subpart \py{"Décimal" if deci12 else "Non Décimal"}
\subpart \py{"Décimal" if deci34 else "Non Décimal"}
\subpart \py{"Décimal" if deci52 else "Non Décimal"}
\end{subparts}
\end{parts}
\end{solution}
\end{questions}

\end{document}